% Abstract. The 2026 Registrar guide sets no word limit; it requires the
% heading, title, author and conferral year above the text (all supplied by the
% abstract environment in yalefmt.tex), double spacing, and page number "iii".
\begin{abstract}

In contemporary machine learning systems, language often mediates interaction between people and modalities. Text-to-image generators, vision-language models, and voice assistants, for example, assume that human intent can be expressed in words before the system interprets or acts on it. Yet verbal expression is serial and limited in bandwidth, and some mental content is difficult to put into words. RepQA, our benchmark for patient education question answering, provides a narrower motivation from the system side: methods can improve readability, but simplified answers may still omit important content. These limitations motivate this dissertation to explore a complementary interface between mind and world that does not require people to articulate their intent first. This problem divides into three questions: (1) how can machine learning uncover neural organization that is not specified in advance, and how can established neural priors in turn be built into machine-learning models; (2) what representation space and structure can best connect brain activity to the world, and how can the necessary supervision be assembled when neural data are scarce; and (3) how can human intention be connected to behavior in the world?

For the first question, we study a two-way exchange between machine learning and neural priors: machine learning can discover neural organization that is not specified in advance, while established neural organization can guide the design of machine-learning models. HyBRiD realizes the first direction by learning outcome-relevant, high-order relationships among brain regions as weighted hyperedges rather than restricting connectivity to pairwise edges; a multi-head drop-bottleneck makes a maximally-informative-minimally-redundant objective tractable, allowing the hyperedges to be learned end-to-end instead of enumerated and pruned in advance. In the other direction, our neuro-informed attention encoder aggregates subject-specific sets of responsive voxels with learnable queries, while its keys contain only voxel position and multi-atlas parcellation embeddings, separating measured fMRI activity from spatial and parcellation priors that can be shared across subjects. FA-EEG likewise uses frame averaging to make the known mirror symmetry between left and right motor imagery exact for any backbone rather than relearning it from scarce trials. This transfer from neuroscience to machine learning also extends beyond models of neural signals: SWA identifies small-world organization in functional networks of language-model attention heads, relates that organization to fluid reasoning, and preserves it when allocating sparsity during pruning.

For the second question, the argument follows alignment-space selection, alignment learning, language-mediated reconstruction, and a more expressive disentangled generator. PRISM I compares candidate alignment spaces and finds that, among those tested, fMRI aligns more closely with language-model semantic representations than with the visual latents conventionally used for reconstruction. MindLLM learns fMRI-to-language alignment through brain instruction tuning, which creates supervision without new scanning by joining fMRI recordings to a decade of existing annotations through their shared stimulus images. The resulting corpus of 980{,}610 conversations spans four families of brain function: perception, memory and knowledge, language and symbols, and complex reasoning. A decoder trained on this corpus answers open-ended questions in a unified format, improves general decoding, and generalizes to a held-out subject without finetuning. PRISM II uses language for visual reconstruction, structuring the representation into objects, attributes, and relations and generating scenes object by object. This object-disentangled space motivates stronger separation between items in the generator itself. \dedit{} extends the generative side beyond object-centric scene construction by learning one prompt per image item and using grouped cross-attention to confine each region to its associated prompt. This enables broader item-level control through text-, mask-, and reference-image-guided editing and object removal; integration with a neural decoder remains future work.

For the third question, the argument progresses from action meaning to goal realization. At the action level, the Mind World Model jointly learns a mind model that encodes EEG as latent actions and a world model that predicts their visual consequences in synchronized egocentric video. This prediction objective, together with cross-modal alignment, connects brain signals to action intention without trial or action labels. The EEG-derived code also turns adjacent video frames into action-conditioned training examples without recorded actions; because it is computed without access to the target frame, it prevents the action channel from simply carrying the future observation. Task completion often calls for specifying the desired outcome rather than each action, motivating the move to goal-directed planning and policies. At this goal level, PHIL and DySta offer complementary interfaces corresponding to the images and text of the semantic bridge. In the dissertation's embodiment framing, PHIL searches toward a graph state representing the outcome depicted in a goal image, using learned node-to-goal distance heuristics. DySta generates actions directly from a language instruction and multi-frame visual history. It separates visual tokens by temporal persistence, retains a single copy and KV cache for static tokens, and refreshes them through a learned recache gate, combining effective task completion with efficient inference. The three models are evaluated independently; connecting decoded representations to these goal interfaces remains future work.

\end{abstract}
